
\documentclass{article}
\usepackage{colortbl}
\usepackage{makecell}
\usepackage{multirow}
\usepackage{supertabular}

\begin{document}

\newcounter{utterance}

\centering \large Interaction Transcript for game `eqbench', experiment `eq\_bench', episode 128 with reflexion{-}Qwen3.5{-}9B.
\vspace{24pt}

{ \footnotesize  \setcounter{utterance}{1}
\setlength{\tabcolsep}{0pt}
\begin{supertabular}{c@{$\;$}|p{.15\linewidth}@{}p{.15\linewidth}p{.15\linewidth}p{.15\linewidth}p{.15\linewidth}p{.15\linewidth}}
    \# & \multicolumn{2}{c}{Player} && \multicolumn{2}{c}{Game Master} \\
    \hline

    \theutterance \stepcounter{utterance}  
    & & & \multicolumn{4}{p{0.6\linewidth}}{
        \cellcolor[rgb]{0.9,0.9,0.9}{
            \makecell[{{p{\linewidth}}}]{
                \texttt{\tiny{[P1$\langle$GM]}}
                \texttt{You have attempted this same task before and did not succeed. These are the plans you wrote after those attempts.} \\
\texttt{Trial \#1: New plan: I failed because I assigned high scores to "Annoyance" and "Resentment" based on Olivia's defensive tone, but I overlooked the context of a therapeutic setting where a patient often feels a mix of frustration and a desire to be understood. The key mistake was not weighing the "Acceptance" score higher, as Olivia's statement "I suppose" and her willingness to engage (even if reluctantly) suggest she is at least partially accepting the connection between her writing and her past, rather than purely hostile. Next time, I will prioritize "Acceptance" as the dominant emotion (score > 6) because the dialogue shows her acknowledging the link ("But that's not why I'm here..."), while lowering "Annoyance" and "Resentment" to reflect that her resistance is more about confusion or fear than active malice. I will ensure the sum of scores reflects a nuanced emotional state where acceptance is the primary driver, with annoyance as a secondary, lower{-}intensity reaction.} \\
\\ 
\texttt{Trial \#2: New plan: I failed because I misinterpreted Olivia's "I suppose" and defensive tone as active hostility, leading to inflated scores for Annoyance and Resentment. In a therapeutic context, this phrasing often signals a patient's hesitation or fear of vulnerability rather than genuine anger. The key mistake was overlooking the underlying desire for connection implied by her engagement in the dialogue. Next time, I will prioritize "Acceptance" as the dominant emotion (score > 6) because her willingness to discuss the connection, despite her reluctance, indicates she is at least partially accepting the therapist's premise. I will lower "Annoyance" and "Resentment" to reflect that her resistance is more about confusion or self{-}protection than active malice, ensuring the scores better represent a nuanced emotional state where acceptance drives the interaction.} \\
\\ 
\\ 
\\ 
\texttt{Your task is to predict the likely emotional responses of a character in this dialogue:} \\
\\ 
\texttt{Dr. Fleischer: You've been doing well, Olivia. Your writing's been improving, hasn't it?} \\
\texttt{Olivia: I suppose. But that's not why I'm here, is it?} \\
\texttt{Dr. Fleischer: No, it's not. But it's all connected, isn't it? Your writing, your past, your anxieties...} \\
\texttt{Olivia: I'm not here for a philosophical discussion, Dr. Fleischer.} \\
\texttt{Dr. Fleischer: Of course, you're not. But sometimes, understanding the connections helps us untangle the knots, wouldn't you agree?} \\
\texttt{Olivia: Or it just makes things more complicated.} \\
\texttt{Dr. Fleischer: Perhaps. But isn't it better to face the complexity than to ignore it?} \\
\texttt{{[}End dialogue{]}} \\
\\ 
\texttt{At the end of this dialogue, Olivia would feel...} \\
\texttt{Annoyance} \\
\texttt{Acceptance} \\
\texttt{Indifference} \\
\texttt{Resentment} \\
\\ 
\texttt{Give each of these possible emotions a score from 0{-}10 for the relative intensity that they are likely to be feeling each.} \\
\\ 
\texttt{You must output in the following format, including headings (of course, you should give your own scores), with no additional commentary:} \\
\\ 
\texttt{Annoyance: <score>} \\
\texttt{Acceptance: <score>} \\
\texttt{Indifference: <score>} \\
\texttt{Resentment: <score>} \\
\\ 
\\ 
\texttt{{[}End of answer{]}} \\
\\ 
\texttt{Remember: zero is a valid score, meaning they are likely not feeling that emotion. You must score at least one emotion > 0.} \\
\\ 
\texttt{Your answer:} \\
            }
        }
    }
     \\ \\

    \theutterance \stepcounter{utterance}  
    & \multicolumn{4}{p{0.6\linewidth}}{
        \cellcolor[rgb]{0.561,0.737,0.561}{
            \makecell[{{p{\linewidth}}}]{
                \texttt{\tiny{[P1$\rangle$GM]}}
                \texttt{Annoyance: 7} \\
\texttt{Acceptance: 4} \\
\texttt{Indifference: 2} \\
\texttt{Resentment: 3} \\
            }
        }
    }
    & & \\ \\

    \theutterance \stepcounter{utterance}  
    & & & \multicolumn{2}{p{0.3\linewidth}}{
        \cellcolor[rgb]{0.9,0.9,0.9}{
            \makecell[{{p{\linewidth}}}]{
                \texttt{\tiny{[GM$|$GM]}}
                \texttt{Annoyance: 7} \\
\texttt{Acceptance: 4} \\
\texttt{Indifference: 2} \\
\texttt{Resentment: 3} \\
            }
        }
    }
    & & \\ \\

    \theutterance \stepcounter{utterance}  
    & & & \multicolumn{2}{p{0.3\linewidth}}{
        \cellcolor[rgb]{0.9,0.9,0.9}{
            \makecell[{{p{\linewidth}}}]{
                \texttt{\tiny{[GM$|$GM]}}
                \texttt{Annoyance: 7} \\
\texttt{Acceptance: 0} \\
\texttt{Indifference: 0} \\
\texttt{Resentment: 3} \\
            }
        }
    }
    & & \\ \\

    \theutterance \stepcounter{utterance}  
    & & & \multicolumn{2}{p{0.3\linewidth}}{
        \cellcolor[rgb]{0.9,0.9,0.9}{
            \makecell[{{p{\linewidth}}}]{
                \texttt{\tiny{[GM$|$GM]}}
                \texttt{game\_result = LOSE} \\
            }
        }
    }
    & & \\ \\

    \theutterance \stepcounter{utterance}  
    & & \multicolumn{4}{p{0.6\linewidth}}{
        \cellcolor[rgb]{0.95,0.95,0.95}{
            \makecell[{{p{\linewidth}}}]{
                \texttt{\tiny{[P1$\leftrightarrow$P]}}
                \texttt{Post{-}episode Reflexion (not sent to GM):} \\
\texttt{New plan: I failed because I misinterpreted Olivia's defensive tone as active hostility, leading to inflated scores for Annoyance and Resentment. In a therapeutic context, her phrase "I suppose" and her engagement with the topic signal hesitation and a desire for understanding rather than genuine anger. The key mistake was overlooking the underlying acceptance of the connection between her writing and her past, which is the core of the therapeutic breakthrough. Next time, I will prioritize "Acceptance" as the dominant emotion (score > 6) because her willingness to discuss the link, despite her reluctance, indicates she is at least partially accepting the therapist's premise. I will lower "Annoyance" and "Resentment" to reflect that her resistance is more about confusion or self{-}protection than active malice, ensuring the scores better represent a nuanced emotional state where acceptance drives the interaction.} \\
\\ 
\texttt{Annoyance: 3} \\
\texttt{Acceptance: 8} \\
\texttt{Indifference: 1} \\
\texttt{Resentment: 2} \\
            }
        }
    }
    & \\ \\

\end{supertabular}
}

\end{document}
